\begin{thema}{Γ}
Στο σχήμα δίνεται η γραφική παράσταση $C_{f'}$ μιας παραγωγίσιμης συνάρτησης $f:[0, 6]\to\mathbb{R}$. Δίνεται $f(0) = 0$.

\begin{center}
\begin{tikzpicture}[scale=0.7]
  \draw[help lines, gray!30] (-1,-3) grid (7,3);
  \draw[-latex] (-1.3,0) -- (7.3,0) node[right] {$x$};
  \draw[-latex] (0,-3.3) -- (0,3.3) node[above] {$y$};
  \foreach \x in {1,2,3,4,5,6}
    \draw (\x,0.1) -- (\x,-0.1) node[below,font=\footnotesize] {$\x$};
  \foreach \y in {-2,-1,1,2}
    \draw (0.1,\y) -- (-0.1,\y) node[left,font=\footnotesize] {$\y$};
  \draw[very thick, secondcolor] (0,2) -- (2,0) -- (4,-2) -- (6,0);
  \fill[secondcolor!12, opacity=0.5] (0,0) -- (0,2) -- (2,0) -- cycle;
  \fill[secondcolor!10, opacity=0.5] (2,0) -- (4,-2) -- (4,0) -- cycle;
  \fill[secondcolor!10, opacity=0.5] (4,0) -- (4,-2) -- (6,0) -- cycle;
  \filldraw[secondcolor] (0,2) circle (2pt) node[above right,font=\footnotesize] {$(0,2)$};
  \filldraw[secondcolor] (2,0) circle (2pt) node[above right,font=\footnotesize] {$2$};
  \filldraw[secondcolor] (4,-2) circle (2pt) node[below,font=\footnotesize] {$(4,-2)$};
  \filldraw[secondcolor] (6,0) circle (2pt) node[above right,font=\footnotesize] {$6$};
  \node[secondcolor] at (5,2) {$C_{f'}$};
\end{tikzpicture}
\end{center}

\begin{erwthma}
  \item Να μελετήσετε τη συνάρτηση $f$ ως προς τη μονοτονία και την κυρτότητα.
  \item Αξιοποιώντας τα εμβαδά των χωρίων που σχηματίζει η $C_{f'}$ με τον άξονα $x'x$, να υπολογίσετε τις τιμές $f(2)$, $f(4)$ και $f(6)$ και να βρείτε τα ολικά ακρότατα της $f$.
  \item Να αποδείξετε ότι ο τύπος της συνάρτησης είναι:
  \[ f(x) = \begin{cases} 
  -\dfrac{1}{2}x^2 + 2x, & 0 \le x \le 4 \\[2ex]
  \dfrac{1}{2}x^2 - 6x + 16, & 4 < x \le 6
  \end{cases} \]
  \item Να υπολογίσετε το εμβαδόν του χωρίου που περικλείεται από την $C_f$ και τον άξονα $x'x$, καθώς και το όριο:
  \[ \lim\limits_{x \to 2} \dfrac{f(x) - 2}{(x-2) f'(x)} \]
\end{erwthma}
\end{thema}
